\documentclass[11pt]{article}

% Basic packages
\usepackage[T1]{fontenc}
\usepackage[utf8]{inputenc}
\usepackage[french]{babel}
\usepackage{lmodern}
\usepackage[a4paper,margin=1in]{geometry}
\usepackage{amsmath,amssymb,amsthm,mathtools}
\usepackage{bm}
\usepackage{enumitem}
\usepackage{hyperref}
\usepackage{xcolor}
\usepackage{graphicx}
\usepackage{booktabs}
\usepackage{microtype}

% Theorem environments
\newtheorem{proposition}{Proposition}
\newtheorem{remark}{Remarque}
\newtheorem{conjecture}{Conjecture}
\newtheorem{question}{Question}

% Shortcuts
\newcommand{\R}{\mathbb{R}}
\newcommand{\Pp}{\mathcal{P}_p}
\newcommand{\eps}{\varepsilon}
\newcommand{\iid}{\mathrel{\overset{\mathrm{i.i.d.}}{\sim}}}
\newcommand{\toP}{\xrightarrow{\mathbb P}}
\newcommand{\toD}{\xrightarrow{\mathcal D}}
\newcommand{\RightarrowW}{\Rightarrow}
\newcommand{\W}{W_p}
\newcommand{\permABC}{\textsc{permABC}}
\newcommand{\WABC}{\textsc{WABC}}

% Macros mirrored from preamble/defs.tex so that news.tex compiles standalone.
\def\obsdata{\mathbf{y}}
\def\simdata{\mathbf{z}}
\def\globparam{\beta}
\def\locparam{\mu}
\def\param{\boldsymbol{\theta}}

\title{Notes sur les nouvelles directions pour \permABC}
\author{Notes de travail pour discussion avec les co-auteurs}
\date{\today}

\begin{document}

\maketitle

\begin{abstract}
Ces notes rassemblent un angle théorique possible pour \permABC, centré sur son lien avec
Wasserstein ABC et sur deux régimes asymptotiques naturels en hiérarchique~:
grand nombre de compartiments observés $K \to \infty$, et grande taille de sur-échantillonnage
$M \to \infty$. L'objectif n'est pas de fixer des résultats définitifs, mais de fournir un cadre
clair pour la discussion, d'isoler des énoncés plausibles, et d'identifier les points à
formaliser en travail futur.
\end{abstract}

\tableofcontents

% =====================================================================
\section{Reviews Biometrika : statut point par point}
\label{sec:biometrika_status}

Légende~:
\textcolor{green!55!black}{\textbf{FAIT}} = traité dans la version actuelle de
\texttt{resubmission.tex}~;
\textcolor{orange!85!black}{\textbf{PARTIEL}} = partiellement traité (modification textuelle
seulement, argument heuristique, ou support numérique manquant)~;
\textcolor{red!70!black}{\textbf{À FAIRE}} = non traité.

\subsection*{Rapporteur 1}
\begin{enumerate}[leftmargin=*]
    \item \textcolor{orange!85!black}{\textbf{PARTIEL}} --- \emph{L'ABC sur un SIR déterministe n'a
    pas de sens~; ajouter de la stochasticité.}
    La Section~3.5 introduit maintenant un bruit multiplicatif log-normal
    $\delta_k^{(t)} = \delta_k \exp(\sigma\xi_k^{(t)})$ avec $\sigma=0,05$, et explique pourquoi
    cela rend l'ABC bien posé sur les sorties simulées plutôt que sur la solution de l'EDO.
    \emph{Manque~:} analyse de sensibilité sur~$\sigma$~; pas de validation sur SIR synthétique.
    \item \textcolor{red!70!black}{\textbf{À FAIRE}} --- \emph{Les modèles stochastiques cassent le
    critère d'appariement en grande dimension.} Non discuté.
    \item \textcolor{green!55!black}{\textbf{FAIT}} --- \emph{Comparer à un ABC-SMC standard avec
    schedule décroissant de~$\varepsilon$.} La Figure~4 compare explicitement
    \textbf{ABC-Vanilla, permABC-Vanilla, ABC-SMC, ABC-PMC, et permABC-SMC} (script
    \texttt{fig4bis\_performance\_comparison\_without\_osum.py}). ABC-SMC et ABC-PMC sans
    permutation échouent dès $K$ modéré alors que permABC-SMC réussit, ce qui isole le gain
    propre aux permutations. \emph{À peaufiner~:} une phrase dans le texte autour de Figure~4
    énonçant explicitement ce constat.
\end{enumerate}

\subsection*{Rapporteur 2}
\begin{enumerate}[leftmargin=*]
    \item \textcolor{red!70!black}{\textbf{À FAIRE}} --- \emph{Élargir l'idée des permutations à
    d'autres domaines LFI.} Non fait.
    \item \textcolor{orange!85!black}{\textbf{PARTIEL}} --- \emph{Plus de justification théorique.}
    Une nouvelle Sec.~2.7 « Perspectives asymptotiques » et un appendice
    (\texttt{asymptotics.tex}) discutent désormais du point de vue WABC au niveau des
    compartiments, de la concentration pour $K\to\infty$ sur
    $\arg\min_\beta W_p(\nu_\beta,\nu_\star)$, et de la non-commutation de la limite $M\to\infty$.
    Les propositions sont encore heuristiques.
    \item \textcolor{red!70!black}{\textbf{À FAIRE}} --- \emph{Commentaire sur l'applicabilité plus
    large des extensions permABC-SMC.} Non ajouté.
    \item \textcolor{red!70!black}{\textbf{À FAIRE}} --- \emph{Comparaison avec LFI neuronal et
    bayésien généralisé sur le Gaussian toy.} Non ajouté.
    \item \textcolor{red!70!black}{\textbf{À FAIRE}} --- \emph{Étendre l'étude de simulation à 3--5
    modèles, dont des cas LFI.} Non ajouté.
    \item \textcolor{red!70!black}{\textbf{À FAIRE}} --- \emph{Pousser l'exemple SIR vers des
    modèles épidémiologiques plus complexes.} Non ajouté.
\end{enumerate}

\subsection*{Rapporteur 3}
\begin{enumerate}[leftmargin=*]
    \item \textcolor{orange!85!black}{\textbf{PARTIEL}} --- \emph{Comparer à WABC~; le tri par
    Hilbert est disponible en multivarié.} De nouveaux paragraphes \texttt{\textbackslash new}
    en Sec.~1.2 positionnent \permABC\ vs \WABC, listent les approximations de Bernton~et~al.
    (Sinkhorn, Hilbert, swap), justifient la LSA exacte dans le régime $K$ modéré / dimension
    non-triviale, et notent qu'en 1D \permABC\ et \WABC\ coïncident.
    \emph{Manque~:} comparaison numérique LSA / Hilbert / Sinkhorn sur le toy.
    \item \textcolor{orange!85!black}{\textbf{PARTIEL}} --- \emph{Théorème~1 limité à
    $\varepsilon\to0$~; envisager $n\to\infty$ ou concentration bayésienne en Wasserstein.} Le
    nouveau matériel asymptotique traite $K\to\infty$ et $M\to\infty$ de manière heuristique.
    Le régime $n\to\infty$ n'est pas traité~; pas de borne TV à $\varepsilon$ fini entre
    $\pi_\varepsilon^*$ et $\pi_\varepsilon$.
    \item \textcolor{red!70!black}{\textbf{À FAIRE}} --- \emph{La condition $\varepsilon<\varepsilon^*$
    est difficile à atteindre en pratique.} Non discuté~; pas de recommandation quand
    $\varepsilon\geq\varepsilon^*$.
    \item \textcolor{orange!85!black}{\textbf{PARTIEL}} --- \emph{Figure~2 contre-intuitive
    (prior plus large que l'approximation a posteriori à $M=K$)~; clarifier.} Aujourd'hui la
    Figure~2 a été refaite sur le Gaussian toy canonique avec $K=10$ et $M\in\{10,\dots,250\}$,
    et la légende a été étendue pour décrire explicitement les deux effets asymptotiques
    (marginale de $\sigma^2$ prior-like, marginale de $\mu_1$ concentrée) quand $M$ croît, en
    lien avec l'équation~\eqref{eq:os_pushforward_limit}. \emph{Manque~:} un court paragraphe
    de lecture de la figure dans le texte~; les hyperparamètres ne sont toujours pas donnés dans
    le texte.
    \item \textcolor{red!70!black}{\textbf{À FAIRE}} --- \emph{Utiliser un toy normal--normal en
    forme close pour discuter le rôle de $M$.} Non ajouté.
    \item \textcolor{red!70!black}{\textbf{À FAIRE}} --- \emph{Contraster avec
    \cite*{bornn2017use} sur le rôle de $M$ en ABC de rejet.} Non ajouté.
    \item \textcolor{red!70!black}{\textbf{À FAIRE}} --- \emph{Clarifier l'initialisation de
    permABC-OS à budget fixe~; les propositions par importance aideraient-elles ?} Non discuté.
    \item \textcolor{red!70!black}{\textbf{À FAIRE}} --- \emph{Comparaison numérique avec la
    distance de Wasserstein.} Non ajouté.
\end{enumerate}

\subsection*{Éditeur associé}
Résume R1--R3. Pas de demande supplémentaire.

\subsection*{Éditeur}
\begin{enumerate}[leftmargin=*]
    \item \textcolor{red!70!black}{\textbf{À FAIRE}} --- \emph{Étendre le Lemme~1 et le Théorème~1
    au cas d'ex\ae quo~; définir $\varepsilon^*$ comme la moitié de la distance minimale entre
    observations distinctes, et montrer que le cardinal des appariements acceptés est alors
    constant.} Non fait.
\end{enumerate}

% =====================================================================
\section{Mock review JASA et autres demandes : statut point par point}
\label{sec:jasa_status}

\subsection*{Préoccupations majeures}
\begin{enumerate}[leftmargin=*]
    \item[\textbf{M1.}] \textcolor{orange!85!black}{\textbf{PARTIEL}} --- \emph{Utiliser des
    métriques de qualité postérieure au-delà de $\varepsilon$ (KL, $W_2$, couverture).} Les
    Figures~4 et~6 reportent désormais la distance Wasserstein-2 tranchée entre le nuage de
    particules et une posterior de référence obtenue par MCMC long sur le Gaussian toy
    (justification et références en Sec.~\ref{sec:news_fig_and_sw2}). \emph{Manque~:} les deux
    pickles \texttt{outliers\_4} restent à retraiter~; idem couverture et KL marginale.
    \item[\textbf{M2.}] \textcolor{green!55!black}{\textbf{FAIT}} --- \emph{Ajouter une baseline
    ABC-SMC standard sans permutations.} La Figure~4 contient \emph{à la fois} ABC-SMC et
    ABC-PMC sans permutation, ainsi que ABC-Vanilla, permABC-Vanilla et permABC-SMC. Les
    baselines sans permutation échouent dès $K$ modéré~: le gain est donc directement
    attribuable aux permutations, pas au framework SMC. \emph{À peaufiner~:} une phrase dans le
    texte autour de Fig~4 rendant ce constat explicite.
    \item[\textbf{M3.}] \textcolor{red!70!black}{\textbf{À FAIRE}} --- \emph{Le Théorème~1 dit peu
    à $\varepsilon$ fini~; donner une borne TV à $\varepsilon$ fini ou atténuer le discours.} La
    phrase « at no cost to theoretical consistency » est encore en Sec.~1.4.
    \item[\textbf{M4.}] \textcolor{orange!85!black}{\textbf{PARTIEL}} --- \emph{Formaliser les
    résultats $M\to\infty$ ou intégrer le matériel asymptotique existant.} Une courte Sec.~2.7
    dans le corps du papier et un appendice complet (\texttt{asymptotics.tex}) sont maintenant
    inclus. Les énoncés restent étiquetés « heuristiques ».
    \item[\textbf{M5.}] \textcolor{orange!85!black}{\textbf{PARTIEL}} --- \emph{Le SIR est une
    preuve de concept, pas une validation.} Seuls le bruit log-normal et le paragraphe de réserve
    sont en place. Pas d'expérience synthétique à vérité terrain connue, pas de sensibilité
    sur~$\sigma$, pas de comparaison à une baseline.
    \item[\textbf{M6.}] \textcolor{red!70!black}{\textbf{À FAIRE}} --- \emph{Comparaison numérique
    avec WABC (Hilbert/Sinkhorn).} Comparaison uniquement textuelle.
    \item[\textbf{M7.}] \textcolor{red!70!black}{\textbf{À FAIRE}} --- \emph{Étendre Lemme~1 /
    Théorème~1 aux ex\ae quo.} (Demandé aussi par l'éditeur Biometrika.)
    \item[\textbf{M8.}] \textcolor{red!70!black}{\textbf{À FAIRE}} --- \emph{Élargir le benchmark
    LFI (méthodes neuronales, bayésien généralisé).} Non fait.
\end{enumerate}

\subsection*{Préoccupations mineures}
\begin{enumerate}[leftmargin=*]
    \item[\textbf{m1.}] \textcolor{red!70!black}{\textbf{À FAIRE}} --- \emph{Paragraphe
    introductif en tête de Section~2.} Non ajouté.
    \item[\textbf{m2.}] \textcolor{red!70!black}{\textbf{À FAIRE}} --- \emph{Donner les
    hyperparamètres du Gaussian toy $(a,b,s,n,K)$.} Non ajouté.
    \item[\textbf{m3.}] \textcolor{orange!85!black}{\textbf{PARTIEL}} --- \emph{Discuter la
    Figure~2 à la lumière de l'analyse $M\to\infty$.} La nouvelle légende le fait explicitement~;
    le paragraphe du corps du texte ne conduit pas encore le lecteur à travers la figure.
    \item[\textbf{m4.}] \textcolor{red!70!black}{\textbf{À FAIRE}} --- \emph{Montrer un cas où
    ABC-Gibbs fonctionne bien.} Non ajouté.
    \item[\textbf{m5.}] \textcolor{gray}{\textbf{N/A}} --- \emph{Retirer le markup de suivi de
    modifications avant soumission.} Gardé intentionnellement pour ce tour~; sera résolu juste
    avant soumission.
    \item[\textbf{m6.}] \textcolor{red!70!black}{\textbf{À FAIRE}} --- \emph{Désambiguïser la
    notation $M$ (jeux de données vs.\ compartiments).} Non renommé.
    \item[\textbf{m7.}] \textcolor{red!70!black}{\textbf{À FAIRE}} --- \emph{Taux de particules
    uniques~: utiliser ou supprimer.} Encore proposé en Sec.~2.1, toujours inutilisé.
    \item[\textbf{m8.}] \textcolor{orange!85!black}{\textbf{PARTIEL}} --- \emph{Benchmarker le
    coût de l'assignment.} Le gain $\mathcal{O}(K^2)$ par swap à chaud de la nouvelle Sec.~2.3
    est textuel~; un benchmark quantitatif vs.\ LSA exact est manquant.
    \item[\textbf{m9--m12.}] Ajoutés lors de la dernière session et tous
    \textcolor{red!70!black}{\textbf{À FAIRE}}~: reporter les hyperparamètres du Gaussian toy,
    clarifier l'initialisation de \permABC-OS à budget fixé, discuter le rôle de $\varepsilon^*$
    en pratique, passer à des modèles épidémiologiques plus complexes.
\end{enumerate}

\subsection*{Autres demandes non couvertes par les listes}
\begin{enumerate}[leftmargin=*]
    \item \textcolor{green!55!black}{\textbf{FAIT}} --- \emph{Rendre explicite le lien entre
    \permABC\ et \WABC\ dans le texte principal.} Nouveaux paragraphes
    \texttt{\textbackslash new} en Sec.~1.2, phrase de clôture dans l'abstract et la conclusion.
    \item \textcolor{green!55!black}{\textbf{FAIT}} --- \emph{Donner un récit asymptotique propre
    pour $M\to\infty$ dans le papier.} Sec.~2.3 pushforward + limite $M\to\infty$ formelle,
    Sec.~2.7 synthèse, appendice \texttt{asymptotics.tex}.
    \item \textcolor{green!55!black}{\textbf{FAIT}} --- \emph{Accélérer l'étape d'assignment en
    SMC.} Sec.~2.3 « Fast assignment via warm-started swaps » (nouvelle sous-section).
    \item \textcolor{green!55!black}{\textbf{FAIT}} --- \emph{Remplacer $\varepsilon$-vs-coût par
    qualité-postérieure-vs-coût dans les benchmarks principaux.} Les Figures~4 et~6 utilisent
    maintenant $\mathrm{SW}_2$.
\end{enumerate}

\section{Motivation et objet de ces notes}

L'objet de cette note courte est de présenter une extension possible de l'histoire \permABC\
actuelle sous une forme qui puisse être discutée avec les co-auteurs avant de décider si une
partie de ce matériel doit entrer dans le papier.

Les idées principales sont les suivantes.
\begin{enumerate}[label=(\arabic*)]
    \item \permABC\ peut être vu comme un analogue hiérarchique de \WABC\ au niveau des
    distributions empiriques des \emph{compartiments}.
    \item Cela suggère un régime asymptotique naturel quand le nombre de compartiments
    $K \to \infty$.
    \item Le pont de sur-échantillonnage $M > K$ a un comportement asymptotique qualitativement
    différent~: pour $M$ grand, la marginale globale tend à redevenir prior-like, tandis que les
    variables locales appariées deviennent des objets de type valeur extrême / meilleur
    appariement, souvent quasi-Dirac.
\end{enumerate}

Le but de ces notes est donc d'organiser ces observations en un récit cohérent et d'expliciter
quelles parties semblent rigoureuses, quelles parties sont heuristiques pour l'instant, et
quelles questions peuvent réalistement être converties en résultats formels.

\section{Résumé asymptotique (condensé de \texttt{asymptotics.tex})}
\label{sec:asympt_short}

On note $\widehat\nu_y^K = \tfrac1K\sum_k \delta_{y_k}$ et $\widehat\nu_z^K = \tfrac1K\sum_k \delta_{z_k}$ les
mesures empiriques des compartiments. Avec un coût induit
$d_K(y,z)^p = \tfrac1K\sum_k \rho(y_k,z_k)^p$, on a l'identité clé
\[
  \min_{\sigma\in S_K} d_K(y,z_\sigma) \;=\; W_p\bigl(\widehat\nu_y^K,\widehat\nu_z^K\bigr).
\]
Donc, au niveau des distributions empiriques de compartiments, \permABC\ \emph{est} WABC appliqué
à la mesure empirique des compartiments~; la différence tient à la projection, qui garde en plus
l'étiquetage des locaux optimalement appariés. Ceci fournit un « dictionnaire WABC » naturel pour
la pseudo-postérieure permutation-invariante $\tilde\pi_\varepsilon$.

\paragraph{Limite $K\to\infty$ (heuristique).} Soit $\nu_\beta$ la loi marginale d'un compartiment sous
le modèle et $\nu_\star$ sa limite sous le vrai mécanisme. Si $\widehat\nu_y^K\Rightarrow\nu_\star$
et $\widehat\nu_z^K\Rightarrow\nu_\beta$, alors l'événement d'acceptation devient
asymptotiquement équivalent à $W_p(\nu_\beta,\nu_\star)\le\varepsilon_K$. Pour $\varepsilon_K\to0$
assez lentement, on attend une concentration de type WABC sur
\[
  \beta^\star \in \arg\min_\beta W_p(\nu_\beta,\nu_\star).
\]
C'est un énoncé sur $\beta$, pas sur le vecteur étiqueté $\mu_{1:K}$~: le transport par
permutations \emph{oublie} les labels des locaux.

\paragraph{Limite $M\to\infty$ à $K$ fixé (heuristique, non-commutation).} Sous
sur-échantillonnage, la probabilité d'acceptation tend vers $1$ quand $M$ croît, donc la marginale
non projetée se relaxe vers le prior~:
\[
  \tilde\pi_\varepsilon^{M}(\mathrm d\beta\mid y)\;\longrightarrow\;\pi_\beta(\mathrm d\beta).
\]
En revanche, les locaux \emph{projetés} se concentrent en mécanisme de valeur extrême sur les
meilleurs candidats parmi $M$~: si $\mu\mapsto c_\beta(y_k,\mu)$ admet un unique minimiseur
$\mu_k^\dagger(\beta,y_k)$, alors $\mu_k^{\star,M}\Rightarrow\delta_{\mu_k^\dagger(\beta,y_k)}$.

\paragraph{Conséquence pratique.} La cible inférentielle est bien à $M=K$. Le sur-échantillonnage
$M>K$ est un \emph{pont algorithmique}~: le sens statistique n'est retrouvé qu'au bout du schedule
$M_0>M_1>\dots>M_T=K$. C'est aussi ce que la nouvelle Figure~2 rend visible.

\paragraph{Questions ouvertes (pour le meeting).}
\begin{itemize}[nosep]
    \item Preuve complète de concentration WABC-style pour $\beta$ sous $K\to\infty$ avec taux
    explicites ?
    \item Régime joint $K\to\infty$, $M\to\infty$ avec $M/K\to c>1$ : que devient-on ?
    \item Caractérisation de la limite des locaux projetés comme carte de transport ou couplage
    de Monge ?
    \item Énoncer l'asymptotique grand-$K$ uniquement sur $\beta$ (laisser les $\mu_{1:K}$ hors du
    théorème principal) ?
\end{itemize}

\section{Mise à jour numérique : nouvelle Figure~2, et passage des Figures~4 et~6 à la Wasserstein-2 tranchée}
\label{sec:news_fig_and_sw2}

Trois modifications numériques ont été apportées au papier pendant la session d'édition du
16~avril~2026. Il vaut la peine de les signaler explicitement, avec leur justification et
leur support bibliographique, pour que la logique survive aux futures éditions du corps du texte.

\subsection{Nouvelle Figure~2 : marginales pushforward sous Over-Sampling}

\paragraph{Ce qui a changé.} La version précédente de la Figure~2 traçait la pseudo-postérieure
Over-Sampling $\pi_\varepsilon^{*M}$ sur un Gaussian toy avec paramètre global $\beta$ et
$M \in \{10,\dots,100\}$. La nouvelle version utilise le Gaussian toy canonique de la
Section~3.2 du papier principal --- variance globale $\sigma^2 \sim \mathrm{IG}(a,b)$, moyennes
locales $\mu_k \sim \mathcal N(0, s^2)$, $K=10$ compartiments --- et fait varier $M$ sur
$\{10, 15, 20, 50, 100, 250\}$. Avec $M=K=10$ (violet) la figure retrouve la pseudo-postérieure
\permABC\ standard~; avec $M=250$ (jaune) elle visualise le régime de la limite pushforward.

\paragraph{Pourquoi c'est important.} La nouvelle figure est le pendant numérique de
l'équation~\eqref{eq:os_pushforward_limit}, introduite pendant la révision JASA. Elle rend
deux effets visibles sur le même panneau~:
\begin{itemize}[nosep]
    \item la marginale sur le paramètre global $\sigma^2$ s'aplatit vers le prior quand $M$
    croît, en ligne avec le résultat de dégénérescence que nous prouvons maintenant
    asymptotiquement~;
    \item la marginale sur le paramètre local $\mu_1$ \emph{se concentre} quand $M$ croît, parce
    que la permutation partielle optimale $\sigma^*$ sélectionne la composante simulée la mieux
    appariée parmi $M$ candidates.
\end{itemize}
Cette séparation nette des comportements asymptotiques global/local est exactement ce que la
Figure~2 devait faire passer, et ce n'était pas lisible sur la version précédente.

\subsection{Passage des Figures~4 et~6 à la Wasserstein-2 tranchée}

\paragraph{Ce qui a changé.} Les Figures~4 et~6 du papier principal reportaient auparavant la
tolérance ABC effective $\varepsilon$ atteinte à un budget de calcul donné. Elles reportent
maintenant la \emph{distance Wasserstein-2 tranchée} $\mathrm{SW}_2$ entre le nuage pondéré de
particules produit par chaque algorithme et une postérieure de référence obtenue par MCMC long
(PyMC) sur le Gaussian toy. Les versions $\varepsilon$ et les doubles axes $N_{\mathrm{sim}}$ /
temps sont déplacés dans l'appendice numérique.

\paragraph{Pourquoi $\varepsilon$ seul ne suffit pas.} C'est le point le plus fort soulevé par
les rapporteurs de la mock-JASA~: la tolérance $\varepsilon$ est une quantité
\emph{algorithmique}, pas une mesure de qualité postérieure. Deux méthodes qui ciblent des
pseudo-postérieures différentes ne peuvent pas être comparées sur $\varepsilon$ seul~; de plus,
le critère d'acceptation de \permABC\ est un infimum sur les permutations, donc un $\varepsilon$
plus petit pour \permABC\ ne se traduit pas mécaniquement par une meilleure approximation
postérieure qu'un concurrent de $\varepsilon$ comparable. Reporter la distance réelle à la
postérieure boucle la boucle.

\paragraph{Pourquoi $\mathrm{SW}_2$ plutôt que les alternatives.}
\begin{itemize}[nosep]
    \item \emph{Divergence de Kullback empirique.} Exige une estimation de densité dans l'espace
    augmenté par les compartiments, souffre d'artefacts de discrétisation en dimension
    $\gtrsim 5$, et est asymétrique (il faut choisir entre $\mathrm{KL}(\pi^{\mathrm{ABC}}\Vert\pi)$
    et $\mathrm{KL}(\pi\Vert\pi^{\mathrm{ABC}})$, les deux versions étant mal conditionnées dès
    que les supports diffèrent légèrement). Sur le toy hiérarchique à $K=20$, l'étape
    d'estimation de densité domine déjà la variabilité de l'estimateur, ce qui masque le signal
    qu'on cherche à mesurer.
    \item \emph{Wasserstein-2 exacte.} En dimension jointe $K+1$ déjà non triviale ($K=20$ donne
    un transport en dimension $21$), et le coût de l'estimation Monte Carlo par Sinkhorn à chaque
    point du benchmark est comparable au coût de l'algorithme ABC lui-même.
    \item \emph{Wasserstein-2 tranchée (notre choix).} Estimateur Monte Carlo en temps linéaire
    en la dimension ambiante, métrise la convergence faible sur $\mathbb R^{d}$, possède des
    garanties statistiques précises, et a été explicitement introduite comme divergence ABC par
    Nadjahi~et~al.~\cite{nadjahi2020approximate}.
\end{itemize}

\paragraph{Support bibliographique pour $\mathrm{SW}_2$.} Définition moderne et barycentres~:
Bonneel~et~al.~\cite{bonneel2015sliced}~; métrisation de la convergence faible et propriétés
statistiques~: Nadjahi~et~al.~\cite{nadjahi2020statistical}~; usage comme divergence ABC~:
Nadjahi~et~al.~\cite{nadjahi2020approximate}. Côté données, l'article de référence reste
Bernton~et~al.~\cite{bernton2019approximate} (WABC) dont notre $\mathrm{SW}_2$ côté postérieur
est le pendant. On utilise $n_{\mathrm{proj}}=200$ projections i.i.d.\ via POT.

\paragraph{Portée honnête.} $\mathrm{SW}_2$ ne détecte pas avec forte probabilité un désalignement
de modes orthogonal à toutes les directions de projection échantillonnées~; en grande dimension
le coût d'échantillonnage des projections grandit. Dans notre benchmark ($d \le K+1 = 21$) ce
n'est pas un problème, mais nous l'énonçons explicitement pour qu'un lecteur adversaire ne puisse
pas nous accuser de survendre la métrique.

\subsection{Bénéfice en présence d'outliers pour Under-Matching}

Le passage à $\mathrm{SW}_2$ permet de quantifier un effet qui n'était qu'anecdotique dans la
version $\varepsilon$. Sur le toy avec $K\in\{10,15,20\}$ et $K_{\mathrm{out}}\in\{0,\dots,5\}$,
une régression linéaire sur $\log(\text{ratio de }\mathrm{SW}_2\text{ de permABC-UM vs.
permABC-SMC})$ contre le nombre d'outliers donne des pentes de $-0,12$ à $-0,16$ par outlier,
avec le ratio à $o=5$ typiquement la moitié du ratio à $o=1$. Cela complète, et est cohérent
avec, la description asymptotique d'Under-Matching en Section~2 du papier.

\begin{thebibliography}{9}
\bibitem{bonneel2015sliced}
N.~Bonneel, J.~Rabin, G.~Peyr\'e and H.~Pfister,
\emph{Sliced and Radon Wasserstein barycenters of measures},
J.\ Math.\ Imaging Vis., 51(1):22--45, 2015.
\bibitem{nadjahi2020statistical}
K.~Nadjahi \emph{et al.},
\emph{Statistical and topological properties of sliced probability divergences},
NeurIPS, 2020.
\bibitem{nadjahi2020approximate}
K.~Nadjahi \emph{et al.},
\emph{Approximate Bayesian computation with the sliced-Wasserstein distance},
ICASSP, 2020.
\bibitem{bernton2019approximate}
E.~Bernton, P.~E.~Jacob, M.~Gerber and C.~P.~Robert,
\emph{Approximate Bayesian computation with the Wasserstein distance},
JRSS~B, 81(2):235--269, 2019.
\end{thebibliography}

\section{Points de discussion pratiques pour les co-auteurs}

À ce stade, les questions les plus utiles pour la discussion semblent être~:
\begin{enumerate}[label=(\arabic*)]
    \item La discussion $K \to \infty$ doit-elle apparaître dans le corps du papier, dans le
    supplément, ou seulement dans des notes internes pour l'instant ?
    \item Voulons-nous énoncer l'argument grand-$K$ uniquement pour le paramètre global $\beta$,
    en laissant les paramètres locaux $\mu_{1:K}$ explicitement hors du théorème principal ?
    \item Voulons-nous inclure une proposition formelle expliquant pourquoi le
    sur-échantillonnage grand-$M$ n'est pas une cible inférentielle mais seulement un pont
    algorithmique ?
    \item Vaudrait-il la peine d'ajouter un petit exemple-jouet où le comportement grand-$M$
    peut être montré analytiquement, par ex.\ un modèle local déterministe ou gaussien ?
    \item Y a-t-il une manière nette de formuler les variables locales projetées comme témoins
    de transport, explications latentes sélectionnées, ou appariements de Monge approchés ?
\end{enumerate}

\section{Prochaines étapes suggérées}

Un enchaînement possible des prochaines étapes~:
\begin{enumerate}[label=(\arabic*)]
    \item Nettoyer la notation pour que la comparaison WABC utilise exactement les mêmes
    symboles que le papier actuel.
    \item Décider quelles parties sont destinées à être du matériel théorème/proposition et
    quelles parties doivent rester en discussion/conjecture.
    \item Construire un toy où le phénomène $M \to \infty$ est explicite.
    \item Si les co-auteurs sont d'accord, transformer la partie grand-$K$ en une proposition
    concise liée à la littérature WABC.
\end{enumerate}

\iffalse
\section{Inventaire des modifications suivies dans \texttt{resubmission.tex}}
\label{sec:tracked_changes}

Cette section liste tous les \texttt{\textbackslash new\{\}} et \texttt{\textbackslash old\{\}}
actuellement présents dans le manuscrit, groupés par fichier. Bleu = \texttt{\textbackslash new\{\}}
(ajouté), rouge barré = \texttt{\textbackslash old\{\}} (suppression en attente). Un bloc est
marqué « aujourd'hui » quand il a été introduit pendant la session d'édition du 16~avril~2026~;
les autres blocs étaient déjà présents dans la version précédente et sont gardés pour
traçabilité.

\paragraph{Préambule.}
\begin{itemize}[nosep]
    \item \texttt{preamble/defs.tex}~: redéfinition de \texttt{\textbackslash new} de
    \texttt{\textbackslash textcolor\{blue\}\{\#1\}} vers \texttt{\{\textbackslash color\{blue\}\#1\}}
    pour que \texttt{\textbackslash new\{\ldots\}} puisse s'étendre sur plusieurs paragraphes
    et contenir des équations en mode display. \textit{(aujourd'hui)}
    \item \texttt{preamble/packages.tex}~: ajout des environnements \texttt{proposition} et
    \texttt{remark}, requis par l'appendice asymptotique. \textit{(aujourd'hui)}
\end{itemize}

\paragraph{Abstract (\texttt{sections/abstract.tex}).}
\begin{itemize}[nosep]
    \item \texttt{\textbackslash new}~: une phrase mentionnant le lien WABC et le comportement
    asymptotique quand $K$ ou $M$ croît. \textit{(aujourd'hui)}
\end{itemize}

\paragraph{Section 1 (\texttt{sections/1\_abc\_to\_perm.tex}).}
\begin{itemize}[nosep]
    \item \texttt{\textbackslash old}~: « This reduces the effective complexity from $K!$
    permutations to a canonical representation\ldots » (l.\,55, version précédente).
    \item \texttt{\textbackslash new}~: court paragraphe-pont rappelant qu'en univarié le tri
    coïncide avec la distance de Wasserstein entre mesures empiriques, donc \permABC\ et
    \WABC\ coïncident dans ce régime. \textit{(aujourd'hui)}
    \item \texttt{\textbackslash new}~: bloc existant positionnant \permABC\ par rapport à
    \WABC\ (global vs.\ hiérarchique, nuage non ordonné vs.\ compartiments avec locaux à
    renommage près).
    \item \texttt{\textbackslash new}~: bloc existant comparant les approximations de Sinkhorn,
    Hilbert et greedy-swap de Bernton~et~al.~; avec une phrase \texttt{\textbackslash new}
    supplémentaire \textit{(aujourd'hui)} renvoyant au schéma de swap à chaud de la Section~2.3.
    \item \texttt{\textbackslash old}~: « As a result, we cannot rely on sorting to minimize
    $d(y,z_\sigma)$ in general\ldots » (l.\,63, version précédente).
\end{itemize}

\paragraph{Section 2 (\texttt{sections/2\_sequential.tex}).}
\begin{itemize}[nosep]
    \item \texttt{\textbackslash new}~: formalisme pushforward pour Over-Sampling
    ($\Psi_{\mathbf{y}}^{(M)}$, représentation pushforward de $\pi_\varepsilon^{*M}$) en
    Section~2.3 (version précédente).
    \item \texttt{\textbackslash old}~: paragraphe informel préexistant sur le comportement
    contrasté des marginales $\beta$ et $\mu$ quand $M$ varie (maintenant remplacé par
    l'analyse grand-$M$ formelle ci-dessous).
    \item \texttt{\textbackslash new}~: analyse grand-$M$ de la pseudo-postérieure non projetée
    ($d_M \to \delta_\beta$ p.s., $\tilde\pi_\varepsilon^M(\beta\mid y) \to \pi(\beta)$), avec
    retour à la Figure~\ref{fig:os} (version précédente, étendue durant la dernière session).
    \item \texttt{\textbackslash new}~: \textbf{nouvelle Sous-section~2.3 \emph{Fast assignment
    via warm-started swaps}}. Introduit la stratégie de warm-start depuis la permutation optimale
    de l'itération précédente, argumente qu'un rafinement par swap à $\mathcal{O}(K^2)$ par
    passe suffit typiquement dans une boucle SMC, et réduit le coût amorti de
    $\mathcal{O}(K^3)$ à $\mathcal{O}(K^2)$. \textit{(aujourd'hui)}
    \item \texttt{\textbackslash new}~: \textbf{nouvelle Sous-section~2.7 \emph{Asymptotic
    perspectives}}. Court résumé dans le corps du texte de (i) le point de vue WABC au niveau des
    compartiments, (ii) la concentration pour $K\to\infty$ sur les minimiseurs de
    $\beta \mapsto W_p(\nu_\beta,\nu_\star)$, (iii) la non-commutation $M\to\infty$, avec un
    pointeur vers le nouvel appendice. \textit{(aujourd'hui)}
    \item Paire \texttt{\textbackslash old}~/~\texttt{\textbackslash new} dans la légende de
    Figure~\ref{fig:os}~: l'ancienne légende (« $M$ décroissant de $150$ à $K=10$ ») est
    remplacée par une légende plus informative liant la figure à l'analyse grand-$M$
    (équation~\eqref{eq:os_pushforward_limit}).
\end{itemize}

\paragraph{Section 3 (\texttt{sections/3\_numerical.tex}).}
\begin{itemize}[nosep]
    \item \texttt{\textbackslash new}~: paragraphe de justification ouvrant le benchmark Gaussian
    toy, introduisant la distance Wasserstein-2 tranchée comme métrique directe de qualité
    postérieure et citant~\cite{bonneel2015sliced,nadjahi2020statistical,nadjahi2020approximate}.
    \textit{(aujourd'hui)}
    \item Paire \texttt{\textbackslash old}~/~\texttt{\textbackslash new} sur le paragraphe
    décrivant la Figure~4~: l'ancienne version reportait le gain d'efficacité en termes de
    $\varepsilon$~; la nouvelle le reporte en termes de $\mathrm{SW}_2$ entre le nuage de
    particules et la postérieure de référence, et renvoie explicitement à la version $\varepsilon$
    en appendice. \textit{(aujourd'hui)}
    \item Paire \texttt{\textbackslash old}~/~\texttt{\textbackslash new} sur
    \texttt{\textbackslash includegraphics} et la légende de Figure~\ref{fig:vanilla_vs_smc}~: le
    fichier figure est maintenant \texttt{fig4\_sw2\_nsim\_seed\_42.pdf} (SW$_2$ vs
    $N_{\mathrm{sim}}$) au lieu de la version $\varepsilon$ précédente. \textit{(aujourd'hui)}
    \item Paire \texttt{\textbackslash old}~/~\texttt{\textbackslash new} sur les trois
    paragraphes décrivant la Figure~6 (benchmark OS/UM avec outliers)~: passage à
    $\mathrm{SW}_2$, et le paragraphe UM est augmenté d'un énoncé quantitatif sur le gain
    log-linéaire en nombre d'outliers (pentes $-0,12$ à $-0,16$ par outlier pour
    $K\in\{10,15,20\}$). \textit{(aujourd'hui)}
    \item Paire \texttt{\textbackslash old}~/~\texttt{\textbackslash new} sur
    \texttt{\textbackslash includegraphics} et la légende de Figure~\ref{fig:osum}~: le fichier
    figure est maintenant \texttt{fig6\_sw2\_time\_seed\_42.pdf} (SW$_2$ vs temps).
    \textit{(aujourd'hui)}
    \item Le paragraphe de bruit stochastique SIR ($\delta_k^{(t)} = \delta_k \exp(\sigma\xi_k^{(t)})$,
    $\sigma=0,05$) introduit lors de la révision précédente reste sans markup.
\end{itemize}

\paragraph{Section 4 (\texttt{sections/4\_conclusion.tex}).}
\begin{itemize}[nosep]
    \item \texttt{\textbackslash new}~: paragraphe de clôture résumant le lien WABC, les limites
    $K\to\infty$ et $M\to\infty$, et la réduction de coût par swap à chaud. \textit{(aujourd'hui)}
\end{itemize}

\paragraph{Appendice (\texttt{sections/5\_appendix.tex}).}
\begin{itemize}[nosep]
    \item \texttt{\textbackslash new}~: inclusion complète de \texttt{sections/asymptotics.tex}
    via \texttt{\textbackslash input}, transformant les notes de travail en une vraie section
    d'appendice « Liens asymptotiques avec Wasserstein ABC ». Nécessite les nouveaux
    environnements \texttt{proposition} et \texttt{remark} ajoutés au préambule.
    \textit{(aujourd'hui)}
\end{itemize}

\paragraph{Éléments ouverts nécessitant une décision accepter/rejeter avant soumission.}
\begin{itemize}[nosep]
    \item Les deux blocs \texttt{\textbackslash old} de Section~1.2 (tri-comme-représentation
    canonique, et « cannot rely on sorting » en multivarié) paraissent redondants une fois le
    matériel WABC/Hilbert en place. Recommandation par défaut~: accepter leur suppression.
    \item Le paragraphe \texttt{\textbackslash old} en Section~2.3 discutant le comportement
    marginal de $\beta$ et $\mu$ avec $M$ est remplacé par l'analyse $M\to\infty$ formelle.
    Recommandation par défaut~: accepter sa suppression.
    \item La légende \texttt{\textbackslash old} de Figure~\ref{fig:os} est strictement plus
    courte et moins informative que le remplacement \texttt{\textbackslash new}. Recommandation
    par défaut~: accepter la nouvelle légende.
    \item Les versions \texttt{\textbackslash old} en $\varepsilon$ des Figures~4 et~6
    (paragraphes et légendes) ne sont gardées que pour la traçabilité. Elles sont strictement
    remplacées par les versions $\mathrm{SW}_2$, qui répondent directement au point des
    rapporteurs JASA sur les métriques de qualité postérieure et qui sont appuyées par les
    références listées en Section~\ref{sec:news_fig_and_sw2}. Recommandation par défaut~:
    accepter les nouvelles versions~; déplacer les figures $\varepsilon$ en appendice comme déjà
    fait pour le récit.
    \item Les phrases \texttt{\textbackslash new} de l'abstract et de la conclusion doivent
    rester et être dé-surlignées avant soumission.
    \item L'enveloppement en gros d'un \texttt{\textbackslash new} autour de l'appendice
    asymptotique est un placeholder~: une fois le contenu validé par les co-auteurs,
    l'enveloppement devrait être retiré pour que l'appendice se lise comme du matériel ordinaire.
\end{itemize}
\fi

\end{document}
